% Partie : théorie de la démonstration
% Chapitre : déduction naturelle
% Section : traduction de N2 vers G2

\documentclass[../../../include/open-logic-section]{subfiles}

\begin{document}

\olfileid[fr]{pt}{nat}{tni}

\olsection{Traduction de \Log{N2} vers \Log{G2}}

\begin{prop}\ollabel{prop:N2-to-G2}
Si $\Log{N2c}\ (\Log{N2i}) \Proves \Gamma \Sequent !A$,
alors $\Log{G2c}\ (\Log{G2i}) + \Cut \Proves
\Gamma' \Sequent \Delta$, où $\Gamma'$ est le
multiensemble de !!{formula}s obtenu en effaçant les
étiquettes de $\Gamma$. On prend $\Delta = \{!A\}$
si $!A$ n'est pas~$\lfalse$, et $\Delta = \emptyset$
sinon. Les variantes entre parenthèses correspondent
respectivement l'une à l'autre.
\end{prop}

\begin{proof}
Procédons par récurrence sur la hauteur de la
!!{proof}~$\delta$ de $\Gamma \Sequent !A$.
Si $\delta$ est l'axiome $x:!A \Sequent !A$,
on prend pour $\pi$ l'axiome correspondant
$!A \Sequent !A$, ou $\lfalse \Sequent \ $
lorsque $!A \ident \lfalse$.

Précisons deux opérations employées dans les cas suivants.
Lorsqu'une traduction de prémisse a un conséquent vide,
on peut rétablir la formule $\lfalse$ à droite par
\RightR{\Weakening} si la règle suivante exige une
prémisse portant cette formule. Inversement, si la
construction donne l'absurde à droite alors que
l'énoncé exige un conséquent vide, une coupure avec
l'axiome $\lfalse \Sequent \ $ retire cette occurrence.
Ces deux opérations conservent exactement le contexte.

Si $\delta$ se termine par une règle d'inférence~$R$,
distinguons les cas.
\begin{enumerate}
  \item Si $R$ est \Intro{\lif}, de formule principale
  $!B \lif !C$, et si $x:!B$ apparaît dans le contexte
  de la prémisse mais pas dans celui de la conclusion,
  $\delta$ se termine par
\[
  \AxiomC{}
  \RightLabel{$\delta_1$}
  \Deduce$x:!B, \Gamma \fCenter !C$
  \RightLabel{\Intro{\lif}}
  \UnaryInf$\Gamma \fCenter !B \lif !C$
  \DisplayProof
  \]
  L'hypothèse de récurrence donne une
  !!{proof}~$\pi_1$ de $!B, \Gamma' \Sequent !C$
  (ou de $!B, \Gamma' \Sequent \ $ si
  $!C \ident \lfalse$). On obtient $\pi$ par
  \RightR{\lif}, après une application de
  \RightR{\Weakening} si $!C \ident \lfalse$.
  \item Si $R$ est \Intro{\lif}, mais que $x:!B$
  n'apparaît pas dans le contexte de la prémisse,
  $\delta$ se termine par
\[
  \AxiomC{}
  \RightLabel{$\delta_1$}
  \Deduce$\Gamma \fCenter !C$
  \RightLabel{\Intro{\lif}}
  \UnaryInf$\Gamma \fCenter !B \lif !C$
  \DisplayProof
  \]
  L'hypothèse de récurrence donne une
  !!{proof}~$\pi_1$ de $\Gamma' \Sequent !C$,
  avec rétablissement de l'absurde à droite si
  nécessaire. On applique \LeftR{\Weakening}
  avec~$!B$, puis \RightR{\lif}, pour obtenir $\pi$.
  \item Si $R$ est \Elim{\lif}, $\delta$ se termine par
\[
  \AxiomC{}
  \RightLabel{$\delta_1$}
  \Deduce$\Gamma_1 \fCenter !B \lif !A$
  \AxiomC{}
  \RightLabel{$\delta_2$}
  \Deduce$\Gamma_2 \fCenter !B$
  \RightLabel{\Elim{\lif}}
  \BinaryInf$\Gamma_1, \Gamma_2 \fCenter !A$
  \DisplayProof
  \]
  où $\Gamma = \Gamma_1 \cup \Gamma_2$.
  L'hypothèse de récurrence donne des !!{proof}s
  $\pi_1$ de $\Gamma_1' \Sequent !B \lif !A$
  et $\pi_2$ de $\Gamma_2' \Sequent !B$,
  après rétablissement de l'absurde à droite dans
  la prémisse mineure si nécessaire.
  On construit $\pi$ comme suit :
\[
  \AxiomC{}
  \RightLabel{$\pi_1$}
  \Deduce$\Gamma_1' \fCenter !B \lif !A$
  \AxiomC{}
  \RightLabel{$\pi_2$}
  \Deduce$\Gamma_2' \fCenter !B$
  \Axiom$!A \fCenter !A$
  \RightLabel{\LeftR{\lif}}
  \BinaryInf$!B \lif !A, \Gamma_2' \fCenter !A$
  \RightLabel{\Cut}
  \BinaryInf$\Gamma_1', \Gamma_2' \fCenter !A$
  \DisplayProof
  \]
  Le multiensemble $\Gamma_1' \cup \Gamma_2'$
  peut différer de $\Gamma' = (\Gamma_1 \cup \Gamma_2)'$.
  Par exemple, si $\Gamma_1$ et $\Gamma_2$ contiennent
  tous deux $x:!C$, l'union additive
  $\Gamma_1' \cup \Gamma_2'$ contient deux copies
  de $!C$, tandis que $(\Gamma_1 \cup \Gamma_2)'$
  n'en contient qu'une pour cette même classe.
  Des applications de \LeftR{\Contraction} retirent
  exactement les copies dues aux classes communes.
  Des étiquettes différentes portant une même formule
  conservent, en revanche, leurs copies distinctes.

  Si la prémisse mineure se termine par l'absurde,
  l'hypothèse de récurrence appliquée à $\delta_2$
  donne une !!{proof} de $\Gamma_2' \Sequent \ $ ;
  l'affaiblissement droit permet alors d'employer
  l'arbre ci-dessus. Si la conclusion finale est
  l'absurde, on la retire par la coupure de normalisation
  du conséquent décrite au début. Ces deux cas sont
  indépendants : la forme de la conclusion ne
  détermine pas celle de la prémisse mineure.
  \item Si $R$ est \FalseCl, $\delta$ se termine par
\[
  \AxiomC{}
  \RightLabel{$\delta_1$}
  \Deduce$x: \lnot !A, \Gamma \fCenter \lfalse$
  \RightLabel{\FalseCl}
  \UnaryInf$\Gamma \fCenter !A$
  \DisplayProof
  \]
  L'hypothèse de récurrence donne, après un affaiblissement
  gauche si la décharge classique est vacante, une
  !!{proof} $\pi_1$ de $\lnot !A, \Gamma' \Sequent \ $.
  On obtient $\pi$ par
\[
  \Axiom$!A \fCenter !A$
  \RightLabel{\RightR{\lnot}}
  \UnaryInf$\fCenter !A, \lnot !A$
  \AxiomC{}
  \RightLabel{$\pi_1$}
  \Deduce$\lnot !A, \Gamma' \fCenter $
  \RightLabel{\Cut}
  \BinaryInf$\Gamma' \fCenter !A$
  \DisplayProof
  \]
  Si la conclusion est l'absurde, on retire à nouveau
  sa formule du conséquent selon la convention fixée.
\end{enumerate}

Voici les autres constructions. L'introduction de la
conjonction, de la disjonction et de l'existentielle
emploie la règle droite G2 correspondante, après
rétablissement éventuel d'une prémisse absurde.
Pour une projection de conjonction, on dérive le
séquent de projection par la règle gauche à partir
de l'axiome de la formule projetée, puis on coupe
avec la preuve traduite de la conjonction.
Pour l'élimination de la disjonction, on rétablit au besoin par
affaiblissement les hypothèses déchargées dans les deux branches.
On applique la règle gauche à conséquent partagé dans \Log{G2i}.
Dans \Log{G2c}, la règle gauche ordinaire est suivie d'une contraction
si le conséquent est non vide. On coupe enfin avec la majeure.

L'introduction de la négation emploie sa règle droite
avec le conséquent vide ; si sa décharge est vacante,
on ajoute d'abord la formule active par affaiblissement
gauche. L'élimination de la
négation emploie sa règle gauche sur la preuve
de la formule positive, puis une coupure avec la
preuve de sa négation. L'élimination intuitionniste
de l'absurde est un affaiblissement droit ; si sa
conclusion est elle-même l'absurde, on garde
directement le conséquent vide.

L'introduction universelle et l'introduction
existentielle emploient les règles droites G2.
L'élimination universelle coupe la preuve de la
majeure avec un séquent d'instanciation obtenu
par la règle gauche universelle et un axiome.
L'élimination existentielle applique la règle
gauche à sa branche mineure traduite, puis
coupe avec la preuve de la majeure ; une décharge
vacante requiert d'abord l'ajout de l'instance
à gauche. Les termes restent clos et la fraîcheur
est conservée par effacement des étiquettes.

Dans chaque construction à plusieurs branches,
on contracte uniquement les répétitions des classes
qui étaient communes aux contextes étiquetés.
Le contexte final est donc exactement le multiensemble
annoncé. Les conventions de conséquent vide et
d'absurde singleton se traitent par les deux
opérations initiales. Seule la traduction de
\FalseCl{} utilise un séquent à plusieurs formules
à droite. La traduction d'une !!{proof} de
\Log{N2i} est donc une !!{proof} de \Log{G2i}
augmenté de la coupure.
\end{proof}
















\end{document}
